% Рисунок из тетради (критерий компактности, шаг 2): Ω ⊂⊂ Ω̃ и шар B(x,ρ).
\begin{figure}[htbp]
\centering
\begin{tikzpicture}[scale=1.0]
  \draw[thick, fill=blue!3]
    plot[smooth cycle, tension=0.8]
      coordinates {(0,0) (2.8,1.1) (5.2,0.3) (5.8,-1.5) (3.6,-2.9) (1.0,-2.5) (-0.8,-1.1)};
  \node at (5.0,-0.2) {$\widetilde{\Omega}$};

  \draw[thick, fill=blue!9]
    plot[smooth cycle, tension=0.85]
      coordinates {(1.0,-0.3) (2.7,0.3) (4.0,-0.5) (3.7,-1.9) (2.0,-2.3) (0.9,-1.4)};
  \node at (1.6,-1.9) {$\Omega$};

  \draw[dashed, red!70!black] (2.9,-1.1) circle (0.6);
  \filldraw (2.9,-1.1) circle (1.2pt);
  \node[anchor=south, font=\footnotesize] at (2.9,-1.02) {$x$};
  \draw[-{Latex[length=1.6mm]}] (2.9,-1.1) -- node[below, font=\tiny]{$\rho$} (3.5,-1.1);
\end{tikzpicture}
\caption{Область $\Omega$ погружается в большую ограниченную $\widetilde{\Omega}$.
При $\rho < \dist(\Omega, \partial\widetilde{\Omega})$ шар $B(x,\rho)$ лежит
в $\widetilde{\Omega}$ для каждого $x \in \Omega$, поэтому вес усреднения полный:
$\int_{\widetilde{\Omega}} \omega_\rho(x-y)\,dy = 1$.}
\label{fig:omega-tilde}
\end{figure}
